\honest{Use the Try Harder, Luke}{Eliezer Yudkowsky}{2008}

I open with a line from John McCarthy: ``When there's a will to fail, obstacles can be found.'' That is my whole thesis. The rest of the post is a joke that repeats it.

As a child I thought Luke Skywalker was a cool Jedi. Watching the films again years later, I found he was a whiny teenager.

Yesterday I looked up the scene in which Yoda says ``Do, or do not. There is no try.'' \nb{Yudkowsky's argument about trying is in the previous day's post, ``Trying to Try.''} This post is the joke that goes with it. \nb{Of the two, this is the one that ended up in the Highlights collection.}

I present an ``outtake'' from the scene, which I made up, in which Mark Hamill argues with George Lucas. Luke raises his ship partway out of the swamp, it sinks back, and the script has him give up. Hamill protests. The ship was already rising, which shows it can be done. Yoda, with eight hundred years of experience, has said it is possible. It is Luke's own spaceship, and without it he is stranded. Hamill asks for a scene of Luke still trying a month later, then for one night of trying, then for five minutes. Lucas refuses each time.

\nb{The joke is funny. It also leaves out how the scene ends. In the film, right after Luke walks away, Yoda lifts the ship out of the swamp himself. Luke says ``I don't believe it,'' and Yoda answers, ``That is why you fail.'' In the film, Luke fails because he does not believe. The post leaves this out and makes the scene a lesson about effort.}

My Lucas has an answer for Hamill: he will not halt the story for five minutes while the ship ``bobs in the swamp like a bathtub toy.'' \nb{It is a good answer. Films leave out the boring parts. A year earlier, in ``The Logical Fallacy of Generalization from Fictional Evidence,'' Yudkowsky had warned against treating stories as evidence about the world, because what a story needs is not what a forecast needs.}

Then comes the punch line, which is also my thesis. Lucas explains: ``People wouldn't try for five minutes before giving up if the fate of humanity were at stake.'' That is a claim about real people. My only evidence for it in this post is the pacing of a film, which is the mistake I warned against a year ago. The scene was about one man's spaceship. I raise the stakes to the fate of humanity only in the last line.
